% HD-G28. GPT-6 Astra Pro /root/g34, inherited effort; 21 November 2026.
\section{HD-G28: the intrinsic kernel of two-dimensional bireflection blocks}
\label{sec:two-block-kernel-scope}

Let \(K\subset\mathbf C\) have characteristic zero, and let
\(N\subset GL_{2r}(K)\), \(r\ge5\), be an actual finite
absolutely irreducible group generated by elements \(x\) with
\(\operatorname{rank}(x-1)\le2\). Suppose its complex module
has a system of \(r\) two-dimensional blocks on which its
permutation image is \(S_r\). Denote the kernel on these blocks
by \(T\). Then \(T\) is a characteristic subgroup of the
abstract group \(N\). In particular, every actual
\(G\subset GL_{2r}(K)\) normalizing \(N\) normalizes the
proper nontrivial unital semisimple coefficient algebra
\[
       K\subsetneq K[T]\subsetneq\operatorname{End}_K(K^{2r}).
                                                               \tag{G28.1}
\]
Here \(K[T]\) means the span of the actual matrices of \(T\).
No descent of the initially chosen complex block projectors is
assumed. The conclusion supplies an intrinsic \(K\)-algebra
even when those particular projectors have not been descended.

For the sixfold source problem, this closes the two-dimensional
block alternative for the normal subgroup generated by all
actual bireflections in coefficient ranks \(10,12,14\).
The exact application to the previously proved coefficient and
Jacobian maps is given in paragraph 6.

\paragraph{1. Actual transpositions and a simultaneous tree change of basis.}
Write \(V_{\mathbf C}=\bigoplus_{j=1}^rV_j\), with
\(\dim V_j=2\). An element permuting a cycle of \(\ell\)
blocks has fixed dimension at most two on their direct sum:
a fixed vector is determined by one block value, which must
also be fixed by the product of the transition maps.
The cycle therefore contributes at least \(2(\ell-1)\)
to the codimension of its fixed space.

Consequently every generating element outside \(T\) permutes
exactly one pair of blocks. Its full codimension budget is
already two. It is the identity on all untouched blocks, and
on the exchanged pair has the literal form
\[
       x|_{V_i\oplus V_j}
          =\begin{pmatrix}0&A\\A^{-1}&0\end{pmatrix}
                                                               \tag{G28.2}
\]
in suitable ordered coordinates. The inverse relation follows
because the fixed space on this pair has dimension two.
In particular \(x^2=1\). An inside generator has nonidentity
components on at most two blocks, since every such component
contributes at least one to its fixed-space codimension.

The graph whose edges are the exchanged pairs of the outside
generators is connected. Otherwise its connected-component sums
would be preserved by all the generators, contradicting the
transitive permutation image. Choose a spanning tree and one
actual swap generator on each of its edges. Fix an isomorphism
\(b_1:\mathbf C^2\to V_1\). For an oriented tree edge from
an already based block \(i\) to a new block \(j\), set
\[
                        b_j=x_{ij}|_{V_i}\,b_i.             \tag{G28.3}
\]
The tree has a unique path to each vertex, so these choices are
consistent. Equation (G28.2) makes the same edge map carry
\(b_j\) back to \(b_i\). Thus, under the actual isomorphism
\[
       B:(\mathbf C^2)^{\oplus r}\longrightarrow V_{\mathbf C},
                   (u_j)\longmapsto\sum_j b_j(u_j),          \tag{G28.4}
\]
all chosen generators are the pure block transpositions.
Tree-edge transpositions generate the full symmetric group:
along the path between two vertices, successive conjugations
of adjacent edge transpositions produce their transposition.
Hence \(N\) contains the literal subgroup of all pure block
permutations. In these coordinates,
\[
                              N=T\rtimes S_r.               \tag{G28.5}
\]
The intersection is trivial because a pure permutation fixing
every block is the identity. The map \(B\) is used only to
prove the group identities; it is not asserted to be defined
over \(K\).

\paragraph{2. The base kernel has generators on at most two blocks.}
Replace every outside generator \(x\), whose permutation is
\((i\ j)\), by \(d_x=x p_{ij}^{-1}\), where \(p_{ij}\)
is the pure transposition in (G28.5). This is an actual element
of \(T\) supported on \(i,j\). Retain each inside generator
unchanged. Let \(\mathcal D\) be this resulting set.
The subgroup generated by all \(S_r\)-conjugates of
\(\mathcal D\) is normalized by the pure \(S_r\), lies
in \(T\), and together with \(S_r\) contains all original
generators. Taking the kernel of its permutation projection
therefore proves
\[
     T=\left\langle p d p^{-1}:p\in S_r,
                                      d\in\mathcal D\right\rangle.
                                                               \tag{G28.6}
\]
Every generator on the right is supported on at most two blocks.

Let \(L\subset GL_2(\mathbf C)\) be any coordinate image
of \(T\). The pure permutations identify all coordinate
images with the same literal finite group \(L\), and
\[
                              T\subset L^r.                 \tag{G28.7}
\]
Moreover \(L\) is irreducible. A common invariant line
\(W\subset\mathbf C^2\) would make \(\bigoplus_jW\)
invariant under both \(T\) and the pure permutations, hence
under \(N\), contradicting absolute irreducibility.

We will use the exact derived-subgroup inclusion
\[
                               (L')^r\subset T.             \tag{G28.8}
\]
To prove it, let \(\mathcal A\) be all coordinate entries of
the generators in (G28.6). They generate \(L\).
For \(a,b\in\mathcal A\), permutations place their chosen
entries in block 1 and their possible other entries in two
distinct partner blocks, 2 and 3. The commutator of the two
resulting elements of \(T\) is exactly
\[
                         \operatorname{diag}([a,b],1,\ldots,1).
                                                               \tag{G28.9}
\]
Indeed on block 2 only the first element acts, on block 3 only
the second acts, and elsewhere both are identities. This also
works when either element originally has only one support block.
The subgroup of entries supported solely in block 1 is normal
in the coordinate image \(L\): conjugate by an element of
\(T\) with the desired first entry. It contains every
\([a,b]\), so it contains their normal closure \(L'\).
Permuting block 1 proves (G28.8); factors on distinct blocks
commute as actual matrices.

\paragraph{3. The solvable coordinate case.}
Write \(\operatorname{Rad}(F)\) for the largest solvable
normal subgroup of a finite group \(F\). It is intrinsic:
the product of two solvable normal subgroups is solvable,
by its quotient and intersection exact sequences, so the
finite product of all of them is the unique largest one.
Every group automorphism preserves it.

For \(r\ge5\), \(S_r\) has no nontrivial solvable normal
subgroup. Its only nontrivial proper normal subgroup is the
nonabelian simple group \(A_r\), and \(S_r\) itself is
nonsolvable. If \(L\) is solvable, (G28.7) makes \(T\)
solvable. Every solvable normal subgroup of \(N\) maps to
a solvable normal subgroup of \(S_r\), hence lies in \(T\).
Thus
\[
                         T=\operatorname{Rad}(N),           \tag{G28.10}
\]
which proves characteristicness in this case.

\paragraph{4. The nonsolvable coordinate case and its exact quotient.}
Let \(Z=L\cap\mathbf C^*I_2\), and apply the actual
projective map \(L\to PGL_2(\mathbf C)\).
The published finite subgroup classification gives only
cyclic groups, dihedral groups, \(A_4,S_4,A_5\).
Since the kernel \(Z\) is abelian, the nonsolvable case
is precisely
\[
                               L/Z\simeq A_5.               \tag{G28.11}
\]
A primary locator is Nguyen--van der Put--Top,
\emph{Algebraic subgroups of \(GL_2(\mathbf C)\)},
Indagationes Mathematicae 19 (2008),287--297,
Theorem2, printed p.288, DOI
\url{https://doi.org/10.1016/S0019-3577(08)80004-3}.
Cohen's full finite \(GL_2\) list in (3.1)--(3.2),
pp.392--394 of the 1976 paper, gives the same conclusion.
No classification of a central extension or Schur multiplier
is required below.

Let \(Q\) be the stable term of the finite derived series
of \(L\). Its image in \(A_5\) remains all of \(A_5\)
at every term because \(A_5\) is perfect. In particular
\(Q\subset L'\) and \(Q\to A_5\) is surjective.
By (G28.8), \(Q^r\subset T\). Define
\[
          Z_0=T\cap Z^r,\qquad
          \pi:T\longrightarrow (L/Z)^r\simeq A_5^r
                                                               \tag{G28.12}
\]
by projectivizing each coordinate. Its kernel is exactly
\(Z_0\), and \(Q^r\subset T\) makes it surjective.
The subgroup \(Z_0\) is abelian and normal in \(N\):
base conjugation fixes scalar coordinates and the pure
permutations merely permute them.

Every solvable normal subgroup of \(N\) lies in \(T\)
by its map to \(S_r\), and its image under \(\pi\) is
a solvable normal subgroup of \(A_5^r\), hence trivial.
For completeness, a normal subgroup of a direct product of
nonabelian simple groups is a product of factors. If its
projection to a factor contains a nonidentity element, taking
commutators with that factor gives a nontrivial normal subgroup
inside that factor, hence the full factor; the remaining
projections are then handled the same way. A product of any
nonempty set of \(A_5\) factors is nonsolvable. Thus
\[
                    \operatorname{Rad}(N)=Z_0.              \tag{G28.13}
\]

The quotient is not merely identified by its composition
factors. Its actual isomorphism is
\[
 \Psi:N/Z_0\longrightarrow A_5^r\rtimes S_r,
 \qquad (t p)Z_0\longmapsto(\pi(t),p).                     \tag{G28.14}
\]
The expression \(tp\) is unique by (G28.5), and pure
permutation conjugation permutes the coordinates in
\(\pi(t)\). This proves multiplicativity. Surjectivity
comes from \(Q^r\) and the pure \(S_r\), and the kernel
is exactly \(Z_0\). The inverse takes a coordinate lift
in \(T\) and the indicated pure permutation; different
lifts differ by \(Z_0\).

The base \(A_5^r\) is the unique minimal nontrivial normal
subgroup of this quotient. It is minimal because normal
subgroups inside it are products of factors and \(S_r\)
permutes the factors transitively. Its centralizer in the
whole quotient is trivial: a centralizing element must fix
every factor, so its permutation component is the identity;
each of its remaining coordinates must then be central in
\(A_5\), hence the identity. Any other minimal normal
subgroup has trivial intersection with the base; their
commutator is contained in that intersection, so it would
centralize the base, a contradiction.

The unique minimal normal subgroup is characteristic.
Combining this fact with (G28.13), the exact intrinsic
description of the original kernel is
\[
 T=\text{the inverse image in }N\text{ of the unique minimal
 normal subgroup of }N/\operatorname{Rad}(N).               \tag{G28.15}
\]
This proves characteristicness also when \(r=5\); no
distinction between an abstract top \(A_5\) and a base
\(A_5\) is assumed without (G28.14).

\paragraph{5. The actual algebra over the original field.}
Return to the original matrices in \(GL_{2r}(K)\).
Characteristicness of \(T\) in \(N\) and normality of
\(N\) in \(G\) give the literal identity
\[
          \rho(g)K[T]\rho(g)^{-1}=K[T]\qquad(g\in G).
                                                               \tag{G28.16}
\]
The image of the semisimple group algebra \(K[T]\) in its
given representation is semisimple, by averaging in
characteristic zero. Here and in (G28.1) the notation denotes
that image. It contains the identity.

Scalar extension injects it into the complex block-diagonal
algebra \(\bigoplus_{j=1}^r\operatorname{End}_{\mathbf C}(V_j)\).
This is a proper subalgebra of the full endomorphism algebra,
so the original \(K\)-algebra is proper. It is not just
\(K\): the coordinate image \(L\) is irreducible of
dimension two, whereas a group of scalar matrices fixes
every line. Thus some element of \(T\) has a nonscalar
coordinate and is a nonscalar matrix on \(V\).
This proves every assertion in (G28.1). The algebra is
constructed from an actual characteristic subgroup, so no
Galois-stability assertion about an arbitrary complex block
system enters the argument.

\paragraph{6. The sixfold source application and its exact scope.}
For a reflection-free source of rank \(10\le m\le14\),
let \(N\) be the subgroup generated by all actual
codimension-two elements of the full coefficient group \(G\).
Conjugation preserves their rank, so \(N\triangleleft G\).
There is at least one such element: all codimensions at least
three would give \(\sum c_j\ge36>2m+6\).
Outside the determinant conclusions already proved from
normalized algebras, \(K[N]=\operatorname{End}_K(V)\).
Indeed it is a nontrivial semisimple normalized algebra;
if proper it is already covered by HD-S38 or HD-S39 together
with the exact smaller-block maps. The nonscalar claim follows
because a nonidentity scalar has rank \(m>2\).
Scalar extension of this equality makes \(N\) absolutely
irreducible: any invariant complex subspace is invariant under
the full endomorphism algebra.

The bireflection specialization of Guralnick--Saxl,
\emph{Generation of finite almost simple groups by conjugates},
Journal of Algebra 268 (2003),519--571,
Theorem8.4, printed pp.557--558, gives in its two-dimensional
block alternative exactly \(m=2r\), \(N/T=S_r\).
Thus \(m=10,12,14\) gives \(r=5,6,7\), and the complete
argument above supplies the proper actual algebra \(K[T]\).
This is precisely the algebra input in HD-S38/HD-S39, not a
replacement of it by complex dimension data.

For Gaussian sources HD-S38 now gives determinant \([-1]\)
or \([-3]\). For the seven fields of HD-G27, a
\(K\)-primitive source is excluded by HD-S39. When actual
\(K\)-blocks occur, their rank is at most seven. The rank-one
character is rational because \(\mu(K)=\{\pm1\}\), rank
two has period rank at most seven, and ranks three through
seven at period rank nine have a reflection by the same
Euler and branch inequalities. HD-G27 gives determinant
\([-1]\) for those faithful blocks.

These deductions retain the actual geometric maps: for a
\(K\)-block \(U\), its stabilizer \(H\), and
\(H_0=\ker(H\to GL(U))\), use
\(Y=\widetilde C/H\) and \(C'=\widetilde C/H_0\).
The coefficient function \(F_v(x)=\pi_U\rho(x)v\), its
inverse \(\sum_{H\backslash G}\rho(x)^{-1}F(x)\), and
the coefficientwise identity \(e_G=e_H\) are those of
HD-G27, paragraphs1--2, or HD-S38, paragraph1.
The actual pullback from \(JC'\) carries the primitive
factor isogenously onto the original factor, with restricted
theta multiplier \(|H_0|\) and determinant norm
\(|H_0|^6\). In the other semisimple-algebra cases, the
literal matrix-unit tensor map and the restriction-of-scalars
field map of HD-S38/HD-S39 are the maps giving their
contradictions. No algebra action on the sixfold is asserted
merely from the existence of the coefficient algebra.

This closes the two-dimensional block alternative. The other
GS alternatives and the passage to still larger coefficient
ranks retain their separately stated proofs and scope.
